\documentclass[a4paper,12pt]{article}
\usepackage[utf8]{inputenc}
\usepackage[T1]{fontenc}
\usepackage[ngerman]{babel}
\usepackage{amsmath}
\usepackage{amssymb}
\usepackage{enumitem}
\usepackage{geometry}
\geometry{a4paper, margin=2.5cm}

\title{Einsendeaufgaben -- Lektion 2\\[0.5em]
\large Modul 61112: Lineare Algebra}
\author{}
\date{}

\begin{document}
\maketitle

\section*{Aufgabe 2.5}

\subsection*{a)}

\textbf{$\forall v \in \mathbb{K}^n: v \sim v$}

Das Nullelement von $\mathrm{GL}_n(\mathbb{K})$ ist $I_n$. Da $H$ eine Untergruppe ist, gilt $I_n \in H$.
Demnach gilt $v \sim v$, da $I_n v = v$.

\bigskip

\textbf{$\forall u, v \in \mathbb{K}^n: u \sim v \Rightarrow v \sim u$}

Sei $u, v \in \mathbb{K}^n$ und $u \sim v$. Es existiert also ein $A \in H$ mit $Au = v$. Da $H$ eine Gruppe mit der Matrixmultiplikation ist, existiert ein $A^{-1} \in H$.
Es gilt also $A^{-1}v = u$ und demnach $v \sim u$.

\bigskip

\textbf{$\forall u, v, w \in \mathbb{K}^n: u \sim v \land v \sim w \Rightarrow u \sim w$}

Sei $u, v, w \in \mathbb{K}^n$ und $u \sim v$ sowie $v \sim w$. Es existiert also ein $A_1$ und $A_2$ mit
\[
  A_1 u = v \quad \text{und} \quad A_2 v = w.
\]
Daraus folgt:
\[
  w = A_2 v = A_2 A_1 u.
\]
Nach der Definition der Verknüpfung einer Gruppe gilt $A_2 A_1 \in H$ und demnach $u \sim w$.

\subsection*{b)}

\textbf{$[0]_{\sim} = \{0\}$}

Es existiert kein $v \in \mathbb{K}^n$ mit $v \neq 0$ und $0 \sim v$, denn es gilt $A0 = 0$ für alle $A \in \mathrm{GL}_n(\mathbb{K})$.
Nach der Symmetrie existiert ebenfalls kein $v \in \mathbb{K}^n$ mit $v \neq 0$ und $v \sim 0$.

\bigskip

\textbf{$[e_1]_{\sim} = \mathbb{K}^n \setminus \{0\}$}

Sei $u, v \in \mathbb{K}^n \setminus \{0\}$. $u$ lässt sich ergänzen zu einer geordneten Basis $B_u$ von $V$, wobei das erste Element $u$ ist.

Wir definieren die Matrix ${}_E M_{B_u}$ durch die geordnete Basis $B_u$. Da ${}_E M_{B_u}$ linear unabhängig ist, existiert eine Inverse ${}_{B_u} M_E$ und demnach
\[
  {}_E M_{B_u} \cdot {}_{B_u} M_E \in \mathrm{GL}_n(\mathbb{K}).
\]
Nach der gleichen Argumentation existieren zwei Matrizen ${}_E M_{B_v},\, {}_{B_v} M_E \in \mathrm{GL}_n(\mathbb{K})$.

Damit können wir $u \sim v$ zeigen:
\[
  {}_E M_{B_v} \cdot {}_{B_u} M_E \cdot u = {}_E M_{B_v} \cdot e_1 = v.
\]

Da bereits $[0]_{\sim} \cup [e_1]_{\sim} = \mathbb{K}^n$, gibt es keine weitere Äquivalenzklasse von $\sim$.

\end{document}
